\section{Introduction}

Continuous glucose monitoring (CGM) generates a dense physiological stream whose forecasted
trajectory informs real-time decisions for people with diabetes: when to eat, when to bolus,
when to adjust activity.
A clinically deployable forecaster must satisfy three properties simultaneously:
(i)~\emph{streaming inference}---constant memory and $\mathcal{O}(1)$ per reading so it can
run on a wearable or mobile device without rebuilding a history window at every step;
(ii)~\emph{probabilistic calibration}---honest quantile intervals that correctly flag
hypoglycemia risk without triggering false alarms; and
(iii)~\emph{directional coherence for plannable actions}---when a user commits to a meal plan,
the forecast should move glucose in the physiologically expected direction.

\textbf{T1DEXI vs.\ AI-READI.}
Shakson Isaac's SSMCGM-Stream \citep{riddell2023t1dexi} was developed on the T1DEXI cohort
(497 participants, all type-1 diabetes), which provides timed insulin bolus logs,
carbohydrate intake, and derived insulin-on-board (IOB) sequences.
This enabled a dose-ordering causal ranking objective that corrects the treatment-by-indication
confound---boluses correlate positively with hyperglycemia in observational data---and produces
genuinely causal counterfactual responses.

AI-READI \citep{aireadi2024} spans a broader phenotypic range (healthy controls,
pre-diabetes lifestyle-controlled, oral/injectable medication, and insulin-dependent
participants) collected at multiple clinical sites.
It provides rich wearable data (CGM, HR, steps, sleep, stress), static clinical variables
(HbA1c, BMI, medications, demographics), and predicted meal flags from a meal-detection model.
However, it does \emph{not} include timed insulin dose or IOB logs as model-input covariates.
\texttt{med\_insulin} records whether a participant uses insulin therapy (a static binary
flag), not when or how much was administered.

This report documents the port of SSMCGM-Stream to AI-READI, covering:
preprocessing and cohort selection (\S\ref{sec:data}),
streaming architecture adaptations (\S\ref{sec:arch}),
training and evaluation protocol (\S\ref{sec:train}),
forecasting results (\S\ref{sec:results}),
personalization diagnostics (\S\ref{sec:pers}),
proxy scenario evaluation (\S\ref{sec:proxy}),
clinical safety metrics (\S\ref{sec:safety}),
interpretability (\S\ref{sec:interp}),
ablations and runtime (\S\ref{sec:ablation}), and
limitations (\S\ref{sec:discuss}).
